\documentclass[10pt,a4paper]{amsart}
\usepackage[margin=19mm]{geometry}
\usepackage{amsmath,amssymb,mathtools}
\usepackage[scr=rsfs,cal=euler]{mathalfa}
\usepackage{xfrac}
\usepackage[unicode,hidelinks,bookmarks=false]{hyperref}
\newcommand{\ie}{i.e.\ }
\newcommand{\cf}{cf.\ }
\newcommand{\resp}{respectively\ }
\newcommand{\Subsection}{\subsection}
\providecommand{\nobreakdash}{\nobreak\hbox{-}}
\setlength{\emergencystretch}{2em}
\multlinegap0pt
\usepackage{bm}
\usepackage{bbm}
\usepackage{dsfont}
\usepackage{xspace}
\usepackage{cancel}
\usepackage{mathtools}
\usepackage{amsthm}
\makeatletter
\@ifundefined{theorem}{\newtheorem{theorem}{Theorem}}{}
\@ifundefined{assumptions}{\newtheorem{assumptions}[theorem]{Assumptions}}{}
\@ifundefined{definition}{\newtheorem{definition}[theorem]{Definition}}{}
\@ifundefined{proposition}{\newtheorem{proposition}[theorem]{Proposition}}{}
\makeatother
\makeatletter
\newcommand{\BL}{\mathbf{BL}}	
\newcommand{\BLN}{\mathbf{BL}}	
\newcommand{\Cu}[1]{\mathbf{C^{#1}}}
\newcommand{\R}{\mathbb{R}}
\newcommand{\TV}{\mathbf{TV}}	
\renewcommand{\d}[1]{\mathrm{d}{#1}}	
\newcommand{\norma}[1]{{\left \|#1\right \|}} 		
\renewcommand{\phi}{\varphi}		
\expandafter\ifx\csname proofof\endcsname\relax
\newenvironment{proofof}[1]
{\smallskip\noindent{\textbf{Proof~of~#1.}}
\hspace{1pt}}{\hspace{-5pt}{\nobreak\nobreak\hfill\nobreak
$\square$\vspace{2pt}\par}\smallskip\goodbreak}
\fi
\newcommand{\reali}{\mathbb{R}}
\newcommand{\ud}{\mathrm{d}}
\makeatother
\title[Asynchronous exponential growth for structured population mo]{Asynchronous exponential growth for structured population models in measure space}
\author{Christian D\"ull, J\'ozsef Z. Farkas, Piotr Gwiazda, Anna Marciniak-Czochra}
\date{Source: arXiv:2607.00658v1 (CC BY 4.0)}
\begin{document}
\maketitle

\noindent\emph{Source and licence.} Asynchronous exponential growth for structured population models in measure space. Version arXiv:2607.00658v1 (CC BY 4.0), distributed under the Creative Commons Attribution 4.0 International licence, as recorded in the arXiv licence metadata for this exact version and shown on the abstract page https://arxiv.org/abs/2607.00658v1. Full rights notice: licenses/arxiv-2607.00658-CC-BY-4.0.txt.\par\smallskip

\noindent\textit{Editorial scope.} Counted unit: the Proposition whose source label is \texttt{prop:abs\_continuity} in the article (source0.tex lines 381--394, source label \texttt{prop:abs\_continuity}), delivered together with the article's own complete original proof of it (source0.tex lines 396--577, one \texttt{proof} environment of 182 lines). No mathematics is written, repaired or re-derived: statement and proof are the article's own text line for line. Required context retained uncounted: Model (lines 212--218); def:WeakSolution (lines 288--305); Formulation:WeakSolution (lines 294--303); Assum (lines 317--328). Every definition, display and label of those lines is retained unchanged. Omitted from this package and not redistributed: 1--211, 219--287, 304--316, 329--380, 395--395, 578--1210. Nothing inside a retained excerpt is omitted or altered: every retained body is delivered exactly as the article's lines are written, so no omission receipt of any kind accompanies this package. Citations: the retained excerpts carry no citation command at all, so no bibliography entry of the article is transcribed and no reference is redistributed. Reference adaptation: none was needed and none was made; every reference inside the delivered unit resolves inside it, so no unresolved reference or citation is produced. Printed slips kept literal: no printed word, symbol, hypothesis or number of the counted statement or of its proof was altered. Layout and class: the article's own class is \texttt{amsart} with packages including amsmath, amsthm, amsfonts, amssymb, mathrsfs, subfigure, graphicx, color, array, booktabs, caption, float, epsfig, latexsym; the wrapper is the lane's pinned \texttt{amsart} editorial wrapper at 10pt on A4, which restates the theorem-like environments the retained text needs and adds the article's own macro definitions that the retained text needs (\texttt{\textbackslash BL}, \texttt{\textbackslash BLN}, \texttt{\textbackslash Cu}, \texttt{\textbackslash R}, \texttt{\textbackslash TV}, \texttt{\textbackslash d}, \texttt{\textbackslash norma}, \texttt{\textbackslash phi}, \texttt{\textbackslash reali}, \texttt{\textbackslash ud}), with extra wrapper packages (bm, bbm, dsfont, xspace, cancel, mathtools). The delivered numbering is the unit's own. Assets: no retained span contains a figure environment, a graphics-inclusion command, a PDF-page inclusion, a table or a TikZ picture; no image file is copied into this package, the package declares no asset dependency and no image file is redistributed with it. Licence: version arXiv:2607.00658v1 of this article is distributed under the Creative Commons Attribution 4.0 International licence, as recorded in the arXiv licence metadata for this exact version and shown on the abstract page https://arxiv.org/abs/2607.00658v1. A full rights notice is delivered at licenses/arxiv-2607.00658-CC-BY-4.0.txt. Characters: every retained excerpt is pure ASCII, so the delivered engine, XeLaTeX with its default Latin Modern text font, reports no missing character.
\begin{align}\label{Model}
\left\{\begin{array}{lll}
\partial_t \mu_t+\partial_x \left(b(x) \, \mu_t \right)+c(x)\,\mu_t &= \displaystyle\int_{\R^+} (\eta(x))(y)\, \ud\mu_t(y),& (t,x) \in  \R^+ \times \R^+,\\ 
    b(0)D_{\lambda}\mu_t(0)&= 0& t \in \R^+,\\
    \mu_{t=0}&=\mu_0,
\end{array}\right.
  \end{align}

\begin{definition}
 \label{def:WeakSolution}
Given $T>0$, a function $\mu_{\bullet} \colon [0,T] \to \mathcal{M}^+(\R^+)$ is a \textbf{measure solution}
of model \eqref{Model}, if $\mu_{\bullet}$ is narrowly continuous, and for all $\phi \in (\Cu{1}
 \cap \BL) \left(\reali^+\times \reali^+\right)$ the following equality holds 
 \begin{align}
 \begin{split}
 \label{Formulation:WeakSolution}
  & \int_{\reali^+} \phi(T,x) \; \d{\mu_T}(x) - \int_{\reali^+}
  \phi(0,x) \; \d{\mu_0}(x)\\
  =&\int_0^T \int_{\reali^+} \left(
   \partial_t \phi(t,x) + b(x) \; \partial_x
   \phi(t,x) - c(x) \; \phi(t,x) \right)\,\ud\mu_t(x)\,\ud{t}   \\
  &+ \int_0^T \int_{\reali^+} \left( \int_{\reali^+} \phi(t,y)
   \d{\left[\eta(x)\right]} (y) \right)\, \ud\mu_t(x)\,\ud{t}.
   \end{split}
\end{align}
\end{definition}

\begin{assumptions}\label{Assum} 
\ \begin{itemize}
\item[(i)] $c \in  \BL(\R^+)\cap C^1(\R^+)$.
\item[(ii)] $x\mapsto \eta(x) \in  \BL\left(\R^+; (\mathcal{M}^+(\R^+), \norma{\cdot}_{\BLN^*})\right)$.
\item[(iii)] $b \in  \BL(\R^+)\cap C^1(\R^+)$, $b>0$. 
\item[(iv)] $b' \leq 0$ .
\item [(v)] There exists a constant $\kappa>0$, such that 
    \begin{align*}
        c(x) \geq \left|c'(x)\right| + \kappa\qquad \forall x \in \R^+.
    \end{align*}
\end{itemize}
\end{assumptions}

\begin{proposition}
\label{prop:abs_continuity}
Let $\mu_{\bullet}$ be a measure solution to \eqref{Model} in the sense of Definition~\ref{def:WeakSolution}.
Then for every
$t\in(0,T]$ there exist $\varepsilon(t)>0$ and $C(t)>0$ such that $\mu_t([0,\varepsilon])\le
C(t)\varepsilon $ 
for all
$0<\varepsilon<\varepsilon(t)$. Consequently,
\[
\limsup_{\varepsilon\to0}
\frac{\mu_t([0,\varepsilon])}{\varepsilon}
<\infty .
\]
\end{proposition}

\begin{proof}
Fix $t>0$. We show that there exist $\varepsilon(t)>0$ and $C(t)>0$ such that
\[
\mu_t([0,\varepsilon])\le C(t)\,\varepsilon
\]
for all $0<\varepsilon<\varepsilon(t)$.


\noindent By Assumptions~\ref{Assum} $b$ is continuous and strictly positive, and thus
for $\delta_1=\frac 1 2 b(0)>0$ there exists $\widetilde\varepsilon_1>0$  with  $b(x)\ge \delta_1$
 for all $x\in[0,\widetilde\varepsilon_1]$. 
Similarly, for $\tilde \varepsilon_2:=\tilde \varepsilon_1+tb(0)$, we can choose $\delta_2>0$ so that 
    \begin{align}
    \label{delta_2_bound}
        b(x)\geq \delta_2\qquad\text{for all }x\in[0,\widetilde\varepsilon_2].
    \end{align}
As $b'\leq 0$ by Assumptions~\ref{Assum}, $\delta_2>\delta_1$. Choose 
\[
\varepsilon(t)
<
\frac12\min\{\widetilde\varepsilon_1,\delta_2 t\}.
\]
\noindent Now let $0<\varepsilon<\varepsilon(t)$ and choose $\psi_\varepsilon\in BL(\R^+)\cap C^1(\R^+)$ satisfying
\begin{align*}
0\le \psi_\varepsilon\le 1,\qquad
\psi_\varepsilon\equiv 1 \ \text{on }[0,\varepsilon],
\qquad
\operatorname{supp}\psi_\varepsilon\subset[0,2\varepsilon].
\end{align*}
Then
\begin{align}
\label{estimate_mu_t}
\mu_t([0,\varepsilon])
\le
\int_{\R^+}\psi_\varepsilon(x)\,\ud\mu_t(x).
\end{align}


\noindent Let $\tilde \varphi_{\varepsilon,t}$ denote the solution of the adjoint
transport-decay problem
\begin{align}
\label{adjusted_dual_problem}
\left\{
\begin{array}{lll}
\partial_{\tau}\tilde \varphi_{\varepsilon,t}
+b\,\partial_x\tilde \varphi_{\varepsilon,t}
-c\,\tilde \varphi_{\varepsilon,t}
=0,
& \text{in } [0,t]\times\R^+,\\[1mm]
\tilde \varphi_{\varepsilon,t}(t,\cdot)
=\psi_\varepsilon,
& \text{in } \R^+.
\end{array}
\right.
\end{align}
It holds that
\[
\tilde \varphi_{\varepsilon,t}\in
C^1([0,t]\times\R^+)\cap BL([0,t]\times\R^+).
\]


\noindent For $\tau\in [0,t]$ the method of characteristics implies
\[
\tilde \varphi_{\varepsilon,t}(\tau,x)
=
\psi_\varepsilon(X_b(t-\tau,x))
\exp\!\left(
-\int_\tau^t
c(X_b(s-\tau,x))\,\ud s
\right),
\]
where $X_b$ denotes the flow of the vector field $b$, i.e. it solves the ODE
	\begin{align}
    \label{flow}
	\partial_s X_b(s,x) = b(X_b(s, x)) \quad \quad X_b(0, x) = x.
	\end{align}
We estimate
\begin{align}
\label{estimate_tilde_varphi}
0\le
\tilde \varphi_{\varepsilon,t}(\tau,x)
\le
e^{t\|c\|_\infty}
\mathbf 1_{\{X_b(t-\tau,x)\in[0,2\varepsilon]\}}.
\end{align}

\noindent
Using \eqref{flow} we compute
\begin{align}
\label{flow_equation}
X_b(t-\tau,x)=x+\int_0^{t-\tau}b(X_b(s,x))\,\ud s.
\end{align}
Now if $x\geq\tilde \varepsilon_1$, then the strict positivity of $b$ together with \eqref{flow_equation} implies 
    \begin{align}
    \label{estimate_above_eps_1}
        X_b(t-\tau,x)\geq x\geq\tilde \varepsilon_1>2\varepsilon\qquad \forall x\in [\tilde\varepsilon_1,\infty).
    \end{align}
To show a similar bound if  $x\in[0,\tilde\varepsilon_1)$, note that  \eqref{flow_equation}  together with $b'\leq 0$ implies 
\begin{align*}
    X_b(s,x)\leq x+sb(0)\leq \tilde \varepsilon_1+tb(0)=\tilde \varepsilon_2\qquad \forall \,0\leq s\leq t.
    \end{align*}
Thus, for all $\tau\leq t-\frac {2\varepsilon}{\delta_2}$ we compute with \eqref{flow_equation} and \eqref{delta_2_bound}
\begin{align}
\label{estimate_below_eps_1}
X_b(t-\tau,x)\ge x+\delta_2(t-\tau)>2\varepsilon\qquad \forall x\in [0,\tilde \varepsilon_1).
\end{align}
Combining  \eqref{estimate_above_eps_1} and \eqref{estimate_below_eps_1} with the estimate \eqref{estimate_tilde_varphi} yields
    \begin{align*} 
\tilde \varphi_{\varepsilon,t}(\tau,\cdot)\equiv0 \qquad \forall \tau\leq t-\frac {2\varepsilon} {\delta_2}.
    \end{align*}
\noindent Using $\tilde \varphi_{\varepsilon,t}$ as a test function in the weak
formulation \eqref{Formulation:WeakSolution} with the upper time boundary replaced by $t$  yields
    \begin{align*}
&\int_{\R^+}\psi_\varepsilon(x)\,\ud\mu_t(x)
=\int_{\R^+}\tilde\varphi_{\varepsilon,t}(t,x)\,\ud\mu_t(x)\\
\iffalse
=&\int_0^t\tilde\varphi_{\varepsilon,t}(0,x)\,\ud\mu_0(x)\\
&+\int_0^t \int_{\reali^+} \left(
   \partial_{\tau}\tilde\varphi_{\varepsilon,t}(\tau,x) + b(x) \; \partial_x
   \tilde\varphi_{\varepsilon,t}(t,x) - c(x) \; \tilde\varphi_{\varepsilon,t}(t,x) \right)\,\ud\mu_{\tau}(x)\,\ud{\tau}   \\
   &+\int_0^t
\int_{\R^+}
\left(
\int_{\R^+}
\tilde \varphi_{\varepsilon,t}(\tau,y)\,
\ud[\eta(x)](y)
\right)
\ud\mu_\tau(x)\,\ud\tau\\
\fi
=&\int_0^t\tilde\varphi_{\varepsilon,t}(0,x)\,\ud\mu_0(x)+
\int_0^t
\int_{\R^+}
\left(
\int_{\R^+}
\tilde \varphi_{\varepsilon,t}(\tau,y)\,
\ud[\eta(x)](y)
\right)
\ud\mu_\tau(x)\,\ud\tau
\end{align*}
as the transport-decay terms cancel
by construction of $\tilde \varphi_{\varepsilon,t}$. Since $\tilde \varphi_{\varepsilon,t}(\tau,\cdot)\equiv0$ for $\tau\leq t-\frac {2\varepsilon}{\delta_2}$ this reduces to
\[
\int_{\R^+}\psi_\varepsilon(x)\,\ud\mu_t(x)
=
\int_{t-\frac{2\varepsilon}{\delta_2}}^t
\int_{\R^+}
\left(
\int_{\R^+}
\tilde \varphi_{\varepsilon,t}(\tau,y)\,
\ud[\eta(x)](y)
\right)
\ud\mu_\tau(x)\,\ud\tau.
\]

\noindent Using the boundedness of $\tilde \varphi_{\varepsilon,t}$ we obtain
\[
\int_{\R^+}\psi_\varepsilon(x)\,\ud\mu_t(x)
\le
e^{t\|c\|_\infty}\|\eta\|_{\BLN}
\int_{t-\frac{2\varepsilon}{\delta_2}}^t
\|\mu_\tau\|_{\TV}\,\ud\tau.
\]
The narrow continuity of $\tau\mapsto\mu_\tau$ on $[0,t]$ gives the bound
\[
R_t:=
\sup_{\tau\in[0,t]}
\|\mu_\tau\|_{\TV}
<\infty.
\]
Together with estimate \eqref{estimate_mu_t} we conclude 
\[
\mu_t([0,\varepsilon])
\le
\int_{\R^+}\psi_\varepsilon(x)\,\ud\mu_t(x)
\le
\frac{2}{\delta_2}
e^{t\|c\|_\infty}
M_\eta R_t\,\varepsilon=:C(t)\varepsilon.
\]

\end{proof}

\end{document}
